\input{preamble.tex}
\begin{document}

\begin{titlepage}
\vspace*{18mm}
{\sffamily\small ПРОЕКТНЫЙ ДОКУМЕНТ · ЧЕРНОВИК 0.3 · 10 НОЯБРЯ 2026}\par
\vspace{8mm}
{\sffamily\bfseries\fontsize{26}{30}\selectfont Тенгварская пишущая\\[2pt] машинка ТПМ-2}\par
\vspace{4mm}
{\large Механическое исполнение на челноке Хаммонда:\\ эскизный проект с расчётами}\par
\vspace{14mm}
\begin{center}
\setlength\tgem{13.5mm}
\begin{tikzpicture}[x=\tgu,y=\tgu]
\input{gen/fig_trace_q.tex}
\end{tikzpicture}
\end{center}
\vspace{6mm}
\noindent\begin{minipage}{0.86\textwidth}\small
\textbf{О чём документ.} ТПМ-2 — ручная машинка без мотора. Все знаки отлиты на одном сменном челноке; клавиша поднимает штифт, челнок поворачивается до него, и бумагу о литеру бьёт молоток с постоянной силой пружины. Продвигающие литеры печатаются от левого края клетки, после удара каретка проходит ширину литеры. Для мёртвого знака весь печатающий узел отходит назад на величину, которую запомнила предыдущая буква, и после удара возвращается на нулевой упор. Оттиск совпадает с оттиском ТПМ-1 знак в знак.

\textbf{Что нового в редакции 0.3.} Выбраны кегль и геометрия челнока; описан выбор позиции штифтами; рассчитаны ход клавиши, ступень кода и допуск на зазоры; оценены силы; посчитана погрешность посадки знаков методом Монте-Карло; добавлены таблица кодовых пластин, спецификация узлов, порядок регулировки, программа испытаний с контрольным текстом и этапы макетирования.
\end{minipage}
\vfill
{\small User · продолжение проектов ТПМ-1 и ТПМ-2 (ред. 0.1, 0.2).}
\end{titlepage}

\tableofcontents
\newpage

%=====================================================================
\section{Выбор конструкции}

\subsection{Критерии}

\begin{description}[style=nextline,leftmargin=0pt]
  \item[Простота в эксплуатации.] Наборщику не нужно знать устройство машинки: удобная раскладка, одно правило порядка ударов, исправление одним нажатием, ровный оттиск при любой силе нажатия, смена шрифта или режима без инструмента.
  \item[Точность механизма.] Положение каждого знака задаётся как можно меньшим числом звеньев; каждое звено доводится жёстким упором или фиксатором до целой единицы; регулировка, от которой зависит посадка всех знаков, одна.
\end{description}

\subsection{Решения}

\begin{center}\small
\begin{tabularx}{\linewidth}{@{}>{\raggedright}p{30mm}>{\raggedright}p{30mm}>{\raggedright\arraybackslash}X@{}}
\toprule
Вопрос & Решение & Почему \\
\midrule
печатающий механизм & челнок Хаммонда & один литероноситель и одна точка удара; молоток бьёт с постоянной силой; челнок сменный \\
как попасть на ось & отвод всего печатающего узла & одно подвижное звено с нулевым упором вместо двух салазок \\
кегль и челнок & em = 1/6'', 4 ряда $\times$ 22 & челнок по размеру близок к хаммондовскому; знаки остаются читаемыми (\S\ref{sec:size}) \\
выбор позиции & штифты, как у Хаммонда & жёсткий упор на каждую позицию \\
ход и ступень кода & $T = 14$~мм, $k = 0{,}6$~мм/u & сброс памяти успевает до первого замера (\S\ref{sec:budget}) \\
нижняя ось & сохранить & без неё подстрочный знак уходит на 0,26~мм, под hyarmen — на 0,8~мм \\
«назад на литеру» & включить & исправление одним нажатием вместо до 16 ударов забоя \\
фиксаторы & на всех секторах и на упоре спуска & прощают неточность зазоров меньше половины ступени \\
контроль хода и очерёдность & сохранить & исключают недожатую клавишу и перекрытие нажатий \\
раскладка & как у ТПМ-1 & ярусы пластины независимы; Shift + гласная даёт долгую гласную на ára \\
\bottomrule
\end{tabularx}
\end{center}

\paragraph{Единица и кегль.} В редакции 0.3 кегль уменьшен до em~=~1/6~дюйма ($\approx$ 4,23~мм; по телу это кегль пики, 12~пунктов). Единица подачи $u = \mathrm{em}/16$ = 1/96~дюйма $\approx$ 0,265~мм. Такая единица уже применялась в пропорциональной машинке: у IBM Selectric Composer литеры занимали от 3 до 9 единиц, а единицу можно было выбрать в 1/72, 1/84 или 1/96~дюйма. Все коды литер в единицах $u$ не меняются; меняются только миллиметры.

%=====================================================================
\section{Предшественники}

\begin{description}[style=nextline,leftmargin=0pt]
  \item[Хаммонд (1880-е).] Литеры отлиты одним челноком, который поворачивается на ободе кольцевой наковальни; нужная литера встаёт перед молотком, и тот под действием пружины бьёт бумагу о литеру. Цель схемы — идеальное выравнивание и равномерный оттиск. Сердце машинки — башенка с полукругом вертикальных штифтов: нажатая клавиша поднимает штифт, челнок поворачивается и останавливается на нём; затем срабатывает молоток, штифт опускается, и челнок возвращается в среднее положение. Цельный челнок занимал дугу около 118° и нёс литеры в три ряда; ряд переключали клавиши, поднимавшие наковальню. Пружина молотка взводилась вместе с возвратом каретки. Челнок делали из твёрдой резины, позже из лёгкого металла. Недостатки — заправка листа и невысокая скорость.
  \item[Varityper.] Наследник Хаммонда с пропорциональной подачей: колесо спуска проходит две, три или четыре единицы в зависимости от нажатой клавиши; спуск связан с молотком, чтобы каретка не трогалась раньше удара.
  \item[Комптометр (1887).] Механизм работает от пальца: клавиша задаёт точку останова сектора и сама его поворачивает. Ход клавиши длинный; поздние модели блокировались при недожатой клавише.
\end{description}

%=====================================================================
\section{Принцип}

\subsection{Геометрия}

Пусть ось молотка при нулевом положении узла находится против точки $P$ на бумаге.

\begin{enumerate}
  \item Продвигающая литера отлита на челноке так, что левый край её клетки приходится на ось молотка; ось техты оказывается в $P + a$.
  \item При отпускании клавиши каретка проходит $w$; против оси молотка теперь $P + w$.
  \item Для надстрочного знака узел отходит назад на $b$: удар в $P + w - b = P + a$.
  \item Для подстрочного знака узел отходит на $c = b - d$: удар в $P + w - c = P + a + d$.
\end{enumerate}

\begin{figure}[h]
\centering
\setlength\tgem{48mm}
\begin{tikzpicture}[x=\tgu,y=\tgu]
\input{gen/fig_geometry.tex}
\end{tikzpicture}
\caption{tinco, акут и две точки снизу. Красным — ось молотка при ударе литеры (1) и после подачи (2). Надстрочный знак (3) бьёт с отводом $b$, подстрочный (4) — с отводом $c$. Размеры в $u$.}
\label{fig:geom}
\end{figure}

\subsection{Код литеры}

\begin{center}\small
\begin{tabular}{@{}clll@{}}
\toprule
& Что задаёт & Диапазон & Куда идёт \\
\midrule
$w$ & подачу каретки после удара & 4--16\,u (12 уровней) & упор спуска, сектор W \\
$b$ & отвод узла для надстрочных & 1--8\,u & сектор B \\
$c = b - d$ & отвод узла для подстрочных & 1--9\,u & сектор C \\
\bottomrule
\end{tabular}
\end{center}

У всех литер $b < w$ и $c < w$: отвод никогда не уводит узел за левый край клетки. Мёртвые знаки несут только признак «над / под». Пробелы задают $w = 4$ или $6$ и обнуляют всю память.

\subsection{Проверка на модели}

Модель ТПМ-2 (\code{simulate\_mech} в \code{tools/gen\_shuttle.py}) прогоняется на обоих примерах (\S\ref{sec:examples}) и на контрольном тексте (\S\ref{sec:tests}) вместе с моделью ТПМ-1. Скрипт требует, чтобы каждая литера и каждый знак легли в одну и ту же точку бумаги, и останавливается при любом расхождении. Расхождений нет.

%=====================================================================
\section{Кегль и геометрия челнока}
\label{sec:size}

Шаг позиции на челноке не может быть меньше ширины молотка, а молоток должен покрывать самую широкую литеру — иначе широкие литеры печатались бы не целиком, а соседние задевались бы. Поэтому шаг — 19\,u (самая широкая клетка 16\,u плюс 3\,u промежутка). Позиций 88 (86 литер и 2 резерва). Дуга челнока принята такой же, как у цельного челнока Хаммонда, — 118°.

\begin{center}\small
\begin{tabular}{@{}lrrcrrrrrrr@{}}
\toprule
em & мм & $u$, мм & ряды & дуга & $R$ & высота & инерция & дуга буквы & точка & отвод \\
\midrule
\tabinput{gen/tab_trade.tex}
\bottomrule
\end{tabular}
\end{center}

\noindent{\small Дуга (развёртка ряда), $R$ (радиус челнока при дуге 118°), высота (ряды по 1,35~em), «дуга буквы» (высота лувы), «точка» (высота точки-техты), «отвод» (наибольший ход ползуна, 9\,u) — в миллиметрах. Инерция — относительно выбранного варианта; челнок считается тонкой полосой, так что инерция растёт как дуга $\times$ высота $\times R^2$.}

\paragraph{Выбор: em = 1/6'', 4 ряда по 22.} Это самый лёгкий челнок из рассмотренных: по инерции он в 5 раз легче варианта редакции 0.2 и вдвое легче трёхрядного того же кегля. Радиус 54~мм — того же порядка, что у Хаммонда, где на 118° помещалось по 30 литер обычного кегля. Цена — мелкий кегль: высота лувы 1,7~мм, точка-техта около 0,5~мм. Это кегль книжного набора, и для металлического челнока такие детали не мелки. Четвёртый ряд требует второй клавиши подъёма наковальни: у Хаммонда их было две (заглавные и цифры), здесь одна поднимает челнок на ряд по половине клавиатуры, другая — на два по регистру.

%=====================================================================
\section{Общая схема}

\begin{figure}[h]
\centering
\begin{tikzpicture}[x=1mm,y=1mm]
\fill[gray!8] (0,70) rectangle (160,76);
\draw[thinmech] (0,70) -- (160,70) (0,76) -- (160,76);
\node[lbl] at (22,73) {лист на каретке};
\draw[arr] (150,80) -- node[lbl,above]{подача $w$} (125,80);
\draw[mech,fill=adv!60,rounded corners=1pt] (60,56) rectangle (100,62);
\node[lbl] at (80,59) {наковальня и челнок};
\draw[mech] (62,56) -- (62,46) -- (98,46) -- (98,56);
\draw[mech] (98,46) -- (110,46) -- (110,86) -- (86,86);
\draw[mech,fill=gray!20] (80,84) rectangle (86,88); \node[lbl,anchor=south] at (83,88.5) {молоток};
\draw[arr] (83,84) -- (83,77);
\node[lbl,anchor=west,align=left] at (111,66) {U-образная\\рама};
\fill[axisred] (83,70) circle (0.9);
\draw[mech,fill=gray!12] (40,40) rectangle (130,44);
\node[lbl,anchor=north] at (85,39.5) {ползун: ход назад до 9\,u (2,4 мм), параллельно листу};
\draw[thinmech] (80,44) -- (80,46);
\draw[mech,fill=ink!70] (132,38) rectangle (135,48); \node[lbl,anchor=west,align=left] at (136,43) {нулевой\\упор};
\draw[mech,fill=sledA!60] (34,46) rectangle (37,50); \draw[mech,fill=sledB!60] (30,38) rectangle (33,42);
\node[lbl,anchor=east,text=sledA] at (29,50) {упор $b$};
\node[lbl,anchor=east,text=sledB] at (29,37) {упор $c$};
\draw[{Latex[length=1.6mm]}-,thick,draw=accent] (44,51) -- ++(14,0) node[lbl,right,text=accent]{отвод};
\draw[mech,fill=white,rounded corners=1pt] (2,6) rectangle (52,22);
\node[lbl,align=center] at (27,14) {сектора памяти B, C, W\\с фиксаторами и защёлками};
\draw[sledA,thick] (20,22) -- (35.5,46); \draw[sledB,thick] (30,22) -- (31.5,38);
\draw[mech,fill=white,rounded corners=1pt] (58,6) rectangle (110,22);
\node[lbl,align=center] at (84,14) {кодовые пластины, мерные\\планки R, W, B, C, О, N; штифты};
\draw[mech,fill=white,rounded corners=1pt] (116,6) rectangle (158,22);
\node[lbl,align=center] at (137,14) {спуск: колесо, собачка,\\переставной упор $w$};
\draw[arr] (58,14) -- (52,14); \draw[arr] (110,14) -- (116,14);
\draw[thinmech] (150,22) -- (150,70);
\node[lbl,anchor=west,align=left] at (151,30) {к рейке\\каретки};
\draw[arr] (84,22) -- (84,40) node[pos=0.5,lbl,right]{О};
\end{tikzpicture}
\caption{ТПМ-2, вид сверху (условно). Челнок с наковальней и молоток сидят на одной U-образной раме, рама — на ползуне. В покое ползун прижат пружиной к нулевому упору; для мёртвого знака планка О отводит его до упора $b$ или $c$, выбор делает признак N. Красная точка — ось молотка.}
\label{fig:top}
\end{figure}

\paragraph{Из чего состоит.} Остов машинки Хаммонда (башенка со штифтами, челнок, наковальня, молоток, каретка с пружинным двигателем, колесо спуска), к которому добавлены ползун с упорами, кодовые пластины, шесть мерных планок, переставной упор спуска, три сектора памяти и блокировки. Полный перечень — в \S\ref{sec:spec}.

%=====================================================================
\section{Узлы}

\subsection{Выбор позиции: штифты}

Под челноком — башенка с полукругом из 22 штифтов, по одному на позицию. Каждая клавиша через свою тягу поднимает штифт своей позиции; клавиши передних рядов вдобавок поднимают наковальню на один ряд, клавиша регистра — на два. Челнок, отведённый в исходное положение, отпускается пружиной, поворачивается и останавливается на поднятом штифте. После удара штифт опускается, и челнок возвращается.

\begin{figure}[h]
\centering
\begin{tikzpicture}[x=1mm,y=1mm]
\def\R{27}
\draw[thinmech,ink!40] (0,0) circle (\R);
\draw[mech,line width=2.4pt,draw=adv!90!ink] (-59:\R) arc (-59:59:\R);
\node[lbl,anchor=west] at (60:{\R+2}) {челнок, дуга 118°};
\foreach \i in {0,...,21} {\pgfmathsetmacro\a{-56.3+\i*5.36}
  \fill[ink!60] (\a:{\R-7}) circle (0.55);}
\pgfmathsetmacro\aa{-56.3+14*5.36}
\fill[axisred] (\aa:{\R-7}) circle (1.0);
\node[lbl,text=axisred,anchor=east] at ($(\aa:{\R-7})+(-1.5,0)$) {поднятый штифт};
\fill (0,0) circle (1);
\draw[thinmech] (0,0) -- (\aa:{\R-7});
\node[lbl,anchor=north] at (0,-2) {ось челнока};
\fill[axisred] (0:{\R+4}) circle (0.9);
\draw[arr] ({\R+12},0) -- ({\R+5.5},0);
\node[lbl,anchor=west] at ({\R+12},0) {молоток (за листом)};
\draw[arr,accent] (-70:{\R+6}) arc (-70:-40:{\R+6});
\node[lbl,text=accent,anchor=north] at (-62:{\R+9}) {пружина поворота};
\node[lbl,anchor=north west,align=left,text width=62mm] at (48,-6) {
22 позиции через 5,36°; шаг по дуге 5,03 мм = 19\,u при $R \approx 54$ мм (на рисунке не в масштабе).\\
Ряд (1--4) задают два подъёма наковальни; позицию в ряду — поднятый штифт.};
\end{tikzpicture}
\caption{Башенка со штифтами (вид сверху, условно). Клавиша поднимает штифт своей позиции; челнок поворачивается до него и ставит литеру против молотка.}
\label{fig:pins}
\end{figure}

\paragraph{Развёртка челнока.} Расположение литер по рядам и позициям — на рис.~\ref{fig:shuttle}. Продвигающая литера отлита так, что против оси молотка встаёт левый край её клетки; мёртвый знак — так, что против оси встаёт его якорь.

\begin{figure}[h]
\centering
\setlength\tgem{6.2mm}
\begin{tikzpicture}[x=\tgu,y=\tgu]
\input{gen/fig_shuttle.tex}
\end{tikzpicture}
\caption{Развёртка челнока: 4 ряда по 22 позиции (не в масштабе; при em = 1/6'' ряд — 111 мм). Серое поле — клетка продвигающей литеры, зелёные и коричневые позиции — надстрочные и подстрочные знаки. Красная линия — то, что встаёт против оси молотка.}
\label{fig:shuttle}
\end{figure}

\subsection{Молоток}

Молоток бьёт с силой своей пружины, а не пальца; поэтому тонкие техты и широкие двойные дуги получают одинаковый удар. Как у Хаммонда, пружина молотка взводится возвратом каретки, а ход клавиши только спускает молоток. Это снимает с клавиши самую большую долю работы (\S\ref{sec:forces}). На строку приходится больше ударов, чем на обычной машинке: на каждую букву — сама буква и в среднем один-два знака. Завода пружины должно хватать с запасом на самую длинную строку.

\subsection{Рама и ползун}

Челнок с наковальней, башенка и молоток сидят на одной U-образной раме, так что при отводе они смещаются вместе и ось удара переносится без перекоса. Рама стоит на ползуне с направляющими, параллельными листу; ход — до 9\,u (2,4~мм).

\subsection{Кодовая пластина и мерные планки}

Под клавишей висит кодовая пластина с выступами над мерными планками — длинными качающимися рамками поперёк клавиатуры. Код задан \emph{зазором}: при полном ходе $T$ выступ с зазором $g$ проходит пустой ход и затем ведёт планку на $s = T - g = k\,v$.

\begin{center}\small
\begin{tabular}{@{}clll@{}}
\toprule
Планка & Кто ведёт & Зазор & Что делает \\
\midrule
R & ярус продвигающей литеры & 0 & поднимает защёлки, сектора уходят в ноль \\
W & ярус продвигающей литеры & $T - k\,w$ & ставит упор спуска и сектор W \\
B & ярус продвигающей литеры & $T - k\,b$ & ставит сектор B (упор $b$) \\
C & ярус продвигающей литеры & $T - k\,c$ & ставит сектор C (упор $c$) \\
О & ярус мёртвого знака & 0 & отводит ползун до выбранного упора \\
N & ярус подстрочного знака & 0 & вводит упор $c$ вместо упора $b$ \\
\bottomrule
\end{tabular}
\end{center}

\begin{figure}[h]
\centering
\begin{tikzpicture}[x=1mm,y=1mm]
\draw[mech,fill=white] (2,50) rectangle (14,54); \node[lbl,anchor=south] at (8,54.5) {клавиша tinco};
\draw[mech] (14,51) -- (128,51); \fill (128,51) circle (1.2); \node[lbl,anchor=south] at (128,52.5) {ось};
\draw[mech,fill=dead!70] (30,51) rectangle (110,45);
\node[lbl] at (70,48) {кодовая пластина (нижний ярус)};
\foreach \x/\g/\nm/\col in {38/0/R/accent,62/8.0/W/ink,82/11.0/B/sledA,102/11.6/C/sledB} {
  \draw[mech,fill=white] (\x-2,45) rectangle (\x+2,{45-16+\g*16/14});
  \fill[\col!35,draw=\col] (\x-5,{45-16-2.4}) rectangle (\x+5,{45-16});
  \node[lbl,text=\col] at (\x,{45-16-6}) {планка \nm};
}
\node[lbl,anchor=north,align=center] at (38,{45-16-9}) {сброс\\$g{=}0$};
\node[lbl,anchor=north,align=center] at (62,{45-16-9}) {$w{=}10$\\$g{=}8{,}0$ мм};
\node[lbl,anchor=north,align=center] at (82,{45-16-9}) {$b{=}5$\\$g{=}11{,}0$ мм};
\node[lbl,anchor=north,align=center] at (102,{45-16-9}) {$c{=}4$\\$g{=}11{,}6$ мм};
\draw[dim] (116,51) -- node[lbl,right,align=left]{$T = 14$ мм\\полный ход} (116,35);
\end{tikzpicture}
\caption{Нижний ярус пластины клавиши tinco при $T = 14$~мм и $k = 0{,}6$~мм/u. Сброс идёт без зазора и потому раньше замера; чем больше значение, тем меньше зазор. Полные значения для всех 88 мест — в приложении~\ref{app:plates}.}
\label{fig:plate}
\end{figure}

\paragraph{Два яруса.} Клавиша регистра поднимает наковальню на два ряда и сдвигает мерные планки, чтобы под выступы встал верхний ярус пластины. Ярусы независимы: на одной клавише может быть мёртвый знак (выступы О и N) и продвигающая литера (выступы R, W, B, C).

\subsection{Спуск}

Колесо спуска связано через шестерню с рейкой каретки; пружинный двигатель каретки всё время стремится его повернуть. Зуб колеса соответствует одной единице хода каретки — 0,265~мм; такую мелкость даёт передаточное число, а не мелкий зуб.

\begin{figure}[h]
\centering
\begin{tikzpicture}[x=1mm,y=1mm]
\def\R{22}
\draw[mech,fill=gray!6] (0,0) circle (\R);
\foreach \i in {0,...,39} {\draw[ink!70] ({90+\i*9}:\R) -- ({90+\i*9+4.5}:{\R+2.2}) -- ({90+(\i+1)*9}:\R);}
\fill (0,0) circle (1.2);
\node[lbl] at (0,-5) {колесо спуска};
\draw[arr] (-8,10) arc (130:60:12) node[lbl,above right]{пружина каретки};
\draw[mech,fill=white] (-1.5,{\R+2.4}) -- (1.5,{\R+2.4}) -- (0,{\R-0.5}) -- cycle;
\draw[mech] (0,{\R+2.4}) -- (0,{\R+12}); \node[lbl,anchor=south] at (0,{\R+12}) {держащая собачка Д};
\def\ang{0}
\draw[thinmech,ink!40] (0,0) -- ({\ang}:{\R+8});
\draw[mech,fill=dead] ({\ang}:{\R-0.5}) -- ({\ang+3}:{\R+3}) -- ({\ang-3}:{\R+3}) -- cycle;
\node[lbl,anchor=west] at ({\ang}:{\R+5}) {переставной упор У};
\draw[{Latex[length=1.8mm]}-{Latex[length=1.8mm]},thick,draw=accent] (90:{\R+5}) arc (90:{\ang+4}:{\R+5});
\node[lbl,text=accent] at (45:{\R+10}) {$w$ зубьев};
\draw[mech,fill=gray!15] (-42,-14) circle (5); \node[lbl,anchor=north] at (-42,-20) {шестерня (на валу колеса)};
\draw[mech,dashed,ink!50] (0,0) -- (-42,-14);
\draw[mech] (-64,-8.6) -- (-22,-8.6);
\foreach \x in {-64,-62,...,-24} {\draw[ink!60] (\x,-8.6) -- (\x+1,-6.8);}
\node[lbl,anchor=south] at (-50,-6.5) {рейка каретки};
\node[lbl,anchor=north west,align=left,text width=56mm] at (30,-12) {
\textbf{При отпускании клавиши:}\\
1. Д выходит из зуба;\\
2. колесо идёт до упора У ($w$ зубьев);\\
3. Д падает в зуб;\\
4. У поднимается и возвращается в ноль.};
\end{tikzpicture}
\caption{Спуск. Планка W отводит упор У на $w$ зубьев вперёд; при отпускании каретка проходит ровно $w$ единиц. Мёртвые знаки упор не трогают.}
\label{fig:esc}
\end{figure}

\subsection{Память: три сектора}

Сектора B, C и W поворачивают планки B, C и W. На дуге каждого сектора — зубчатый фиксатор и защёлка. Сектор B держит упор отвода для надстрочных знаков, C — для подстрочных, W — упор возврата «назад на литеру».

\begin{figure}[h]
\centering
\begin{tikzpicture}[x=1mm,y=1mm]
\fill (0,0) circle (1.2);
\draw[mech,fill=sledA!10] (0,0) -- (60:30) arc (60:120:30) -- cycle;
\foreach \i in {0,...,9} {\pgfmathsetmacro\a{114-\i*6}
  \draw[ink!70] (\a:30) -- (\a:32);
  \node[lbl,font=\sffamily\tiny] at (\a:35) {\i};}
\draw[mech,fill=white] (0,40) -- (2,32.5) -- (-2,32.5) -- cycle;
\draw[mech] (0,40) -- (0,46); \node[lbl,anchor=south] at (0,46) {защёлка};
\draw[mech] (0,0) -- (-30,-8);
\node[lbl,anchor=east] at (-30,-8) {от планки B};
\draw[mech,sledA,thick] (0,0) -- (40,-6) -- (70,-6);
\draw[mech,fill=sledA!60] (70,-9) rectangle (73,-3);
\node[lbl,anchor=west,text=sledA] at (74,-6) {упор $b$ ползуна};
\draw[arr] (-20,14) arc (150:125:22); \node[lbl,anchor=east] at (-20,14) {пружина к нулю};
\node[lbl,anchor=north west,align=left,text width=72mm] at (-35,-14) {
\textbf{На каждой продвигающей литере:}\\
1. планка R поднимает защёлку, пружина возвращает сектор в ноль;\\
2. планка B ведёт сектор на $b$ делений;\\
3. защёлка падает в ближайшую впадину и держит сектор.\\
\textbf{На мёртвом знаке} сектор не трогают: он служит упором.};
\end{tikzpicture}
\caption{Сектор памяти B. Сектора C и W устроены так же.}
\label{fig:mem}
\end{figure}

\subsection{Отвод печатающего узла}

При мёртвом знаке планка О через пружинную тягу толкает ползун назад до упора $b$, а если ярус знака несёт выступ N, — до упора $c$. Тяга пружинная, поэтому положение задаёт упор, а не глубина нажатия. После удара пружина возвращает ползун на нулевой упор. В звеньях, задающих положение знака, — нулевой упор и упор сектора; нулевой упор регулируется один раз и переносит посадку всех знаков одинаково.

\subsection{Служебные клавиши}

\begin{itemize}
  \item \textbf{Пробелы 4u и 6u} — постоянные пластины: $w = 4$ или $6$, память в ноль.
  \item \textbf{Забой} — колесо спуска назад на один зуб (1\,u). Память не трогает.
  \item \textbf{Назад на литеру} — храповой толкатель поворачивает колесо спуска назад до упора сектора W и сбрасывает память. Если затем ударить ту же литеру через корректирующую ленту, она ляжет на прежнее место, а память восстановится.
  \item \textbf{$\pm u$} — качающийся рычажок: пока он отклонён, нулевой упор ползуна смещён на единицу. Ручная поправка для редких случаев.
\end{itemize}

\subsection{Блокировки}

\begin{itemize}
  \item \textbf{Контроль полного хода.} Начатая клавиша должна дойти до конца: храповая гребёнка не даёт ей вернуться с полпути.
  \item \textbf{Очерёдность.} Пока клавиша не отпущена — подача не прошла, ползун не вернулся, защёлки не легли, — остальные клавиши заперты.
\end{itemize}

%=====================================================================
\section{Ход клавиши}
\label{sec:budget}

\subsection{Бюджет хода}

Сброс памяти (планка R) должен закончиться раньше, чем любой выступ замера коснётся своей планки, иначе сектор, ещё не сброшенный в ноль, начнут вести дальше. Раньше всех начинает выступ W самой широкой литеры: его зазор $T - 16k$. Принято, что сброс занимает первые 3,5~мм хода, и нужен запас 0,5~мм.

\begin{center}\small
\begin{tabular}{@{}rrrrrrc@{}}
\toprule
$T$, мм & $k$, мм/u & зазор W мин. & зазор B мин. & зазор C мин. & половина ступени & сброс успевает \\
\midrule
\tabinput{gen/tab_budget.tex}
\bottomrule
\end{tabular}
\end{center}

\paragraph{Выбор: $T = 14$~мм, $k = 0{,}6$~мм/u.} Из годных вариантов это самый короткий ход с ещё приемлемой ступенью. $T = 12$~мм проходит только при $k = 0{,}5$, а тогда половина ступени — всего 0,25~мм, и требования к зазорам жёстче (\S\ref{sec:gaptol}). $T = 16$~мм даёт большую ступень, но ход становится утомительно длинным.

\subsection{Порядок действий}

\begin{figure}[h]
\centering
\begin{tikzpicture}[x=7.4mm,y=1mm]
\draw[ink!50,->] (0,0) -- (16.2,0);
\foreach \x in {0,2,...,14} {\draw[ink!50] (\x,0) -- (\x,-1.5) node[lbl,below]{\x};}
\node[lbl,anchor=north] at (7,-5) {ход клавиши, мм};
\node[lbl,anchor=north] at (15.3,-5) {отпуск.};
\def\row#1#2{\node[lbl,anchor=east] at (-0.2,#1) {#2};}
\node[lbl,anchor=west,font=\sffamily\footnotesize\bfseries] at (0,74) {мёртвый знак};
\row{68}{отвод ползуна (О), тяга} \fill[sledA!40] (0,66) rectangle (4,70);
\row{61}{подъём штифта, поворот челнока} \fill[accent!30] (0.5,59) rectangle (10,63);
\row{54}{спуск молотка} \fill[axisred] (13.7,52) rectangle (14,56);
\row{47}{возврат ползуна} \fill[sledA!25] (14.1,45) rectangle (15.8,49);
\node[lbl,anchor=west,font=\sffamily\footnotesize\bfseries] at (0,40) {продвигающая литера};
\row{34}{сброс памяти (R)} \fill[accent!35] (0,32) rectangle (3.5,36);
\row{27}{замер W (с 4,4 мм), B, C (с 8,6 мм)} \fill[sledB!25] (4.4,25) rectangle (14,29); \fill[sledB!50] (8.6,25) rectangle (14,29);
\row{20}{подъём штифта, поворот челнока} \fill[accent!30] (0.5,18) rectangle (10,22);
\row{13}{спуск молотка} \fill[axisred] (13.7,11) rectangle (14,15);
\row{6}{подача $w$; защёлки держат} \fill[axisblue!45] (14.1,4) rectangle (15.8,8);
\end{tikzpicture}
\caption{Порядок действий по ходу клавиши при $T = 14$~мм и $k = 0{,}6$~мм/u. Самое раннее начало замера (W самой широкой литеры) — 4,4~мм; B и C начинают не раньше 8,6~мм. Челнок встаёт на штифт раньше, чем спускается молоток.}
\label{fig:stroke}
\end{figure}

\subsection{Допуск на зазоры кода}
\label{sec:gaptol}

Фиксатор доводит значение до ближайшей единицы, поэтому ошибка зазора меньше половины ступени (0,3~мм) не вредит. Ошибка больше — и литера получит соседнее значение. Если ошибку зазора (изготовление пластины плюс зазоры в шарнирах планки) считать нормальной со стандартным отклонением $\sigma$, вероятность промаха на одном замере — $2\,(1-\Phi(0{,}3/\sigma))$.

\begin{center}\small
\begin{tabular}{@{}rrrrr@{}}
\toprule
$\sigma$, мм & $\pm3\sigma$, мм & запас $0{,}3/\sigma$ & промах на замер & промахов на страницу \\
\midrule
\tabinput{gen/tab_gap.tex}
\bottomrule
\end{tabular}
\end{center}

\noindent{\small Страница — около 1500 продвигающих литер по три замера (допущение).}

\paragraph{Требование.} $\sigma \le 0{,}05$~мм (допуск около $\pm 0{,}15$~мм) — ни одного промаха на десятки тысяч страниц. При $\sigma = 0{,}06$~мм промах уже бывает примерно раз на 400 страниц, при 0,1~мм — десяток на страницу. Это главное требование к точности изготовления.

%=====================================================================
\section{Силы и энергия}
\label{sec:forces}

Молоток взводится возвратом каретки, поэтому клавиша выполняет только замер, подъём штифта и спуски. Оценка для худшего случая ($w = 16$, $b = 8$, $c = 9$) при принятых силах пружин:

\begin{center}\small
\begin{tabular}{@{}lrrr@{}}
\toprule
Работа на ходу продвигающей клавиши & сила, Н & путь, мм & работа, мДж \\
\midrule
планка R (сброс) & 0,3 & 3,5 & 1,05 \\
планка W & 0,3 & 9,6 & 2,88 \\
планка B & 0,3 & 4,8 & 1,44 \\
планка C & 0,3 & 5,4 & 1,62 \\
подъём штифта & 0,5 & 3,0 & 1,50 \\
спуски челнока и молотка, блокировки, трение & & & 1,00 \\
\midrule
\textbf{итого} & & & \textbf{9,5} \\
\bottomrule
\end{tabular}
\end{center}

На ходе 14~мм это в среднем 0,68~Н, около 70~гс. Мёртвая клавиша: отвод ползуна пружинной тягой 1,2~Н на первых 4~мм (4,8~мДж), штифт (1,5~мДж) и спуски (1~мДж) — около 7~мДж, в среднем около 55~гс.

Для сравнения: у рычажной машинки клавишный рычаг с отношением около 4:1 превращает нажатие в 60~гс в 240~гс тяги на пятке рычага. Порядок сил тот же, но в ТПМ-2 сила не зависит от того, насколько сильно бьёт палец: удар задан пружиной молотка.

\paragraph{Чего не хватает.} Силы пружин планок и трение — допущения; их нужно заменить замерами на макете (\S\ref{sec:stages}). Можно частично уравновесить планки встречными пружинами, но тогда хуже возврат.

%=====================================================================
\section{Литеры}

\begin{figure}[h]
\centering
\setlength\tgem{26mm}
\begin{tikzpicture}[x=\tgu,y=\tgu]
\input{gen/fig_slugs_adv.tex}
\end{tikzpicture}
\caption{Продвигающие литеры (увеличено). Красная линия — ось молотка (левый край клетки), серая — ось техты. Под литерой её код.}
\end{figure}

\begin{figure}[h]
\centering
\setlength\tgem{24mm}
\begin{tikzpicture}[x=\tgu,y=\tgu]
\input{gen/fig_slugs_dead.tex}
\end{tikzpicture}
\caption{Мёртвые знаки (увеличено). Якорь знака — на оси молотка; под знаком указано, на какую величину уходит узел.}
\end{figure}

Набор знаков тот же, что в ТПМ-1: один комплект гласных на все тенгвы, черта удвоения и носовая черта в узком и широком вариантах, особая черта для lambe и alda, литеры «ára + гласная».

%=====================================================================
\section{Раскладка}

\begin{figure}[h]
\centering
\begin{tikzpicture}[x=1mm,y=1mm]
\input{gen/keyboard.tex}
\end{tikzpicture}
\caption{Раскладка ТПМ-2 (совпадает с ТПМ-1). Зелёные клавиши — надстрочные знаки (отвод $b$), коричневые — подстрочные (отвод $c$). Вверху фиолетовым — верхний регистр.}
\label{fig:kbd}
\end{figure}

%=====================================================================
\section{Набор: два примера}
\label{sec:examples}

В таблицах: «удар» — точка бумаги, куда пришёлся удар; «отвод» — смещение узла; «каретка» — положение после клавиши; «память» — сектора B / C / W. Все величины — в $u$.

\subsection{Квенья: \emph{Elen síla lúmenn' omentielvo}}

\begin{figure}[h]
\centering
\setlength\tgem{12.5mm}
\begin{tikzpicture}[x=\tgu,y=\tgu]
\input{gen/fig_trace_q.tex}
\end{tikzpicture}
\caption{Квенийский режим. Риски: красные — удары литер, зелёные и коричневые — удары знаков.}
\end{figure}

{\small\setlength\tgem{5mm}
\begin{longtable}{@{}rlrrrc@{}}
\toprule
\# & Клавиша & Удар & Отвод & Каретка & Память $b/c/w$ \\
\midrule\endhead
\tabinput{gen/tab_trace_q.tex}
\bottomrule
\end{longtable}}

\clearpage
\subsection{Синдарин: \emph{pedo mellon a minno}}

\begin{figure}[h]
\centering
\setlength\tgem{14mm}
\begin{tikzpicture}[x=\tgu,y=\tgu]
\input{gen/fig_trace_s.tex}
\end{tikzpicture}
\caption{Синдаринский режим. В «mellon» черта для lambe бьёт с отводом $c = 5$.}
\end{figure}

{\small\setlength\tgem{5mm}
\begin{longtable}{@{}rlrrrc@{}}
\toprule
\# & Клавиша & Удар & Отвод & Каретка & Память $b/c/w$ \\
\midrule\endhead
\tabinput{gen/tab_trace_s.tex}
\bottomrule
\end{longtable}}

%=====================================================================
\section{Точность}

\subsection{Квантование}

Литеры отлиты по точным якорям шрифта, поэтому знак встаёт на ось буквы без ошибки квантования; округление до $u$ затрагивает только промежутки: не больше 0,5\,u на сторону (0,13~мм) и не больше 0,4\,u (0,11~мм) на нижней оси.

\subsection{Механика: Монте-Карло}

Погрешность посадки знака относительно оси буквы складывается из ошибок звеньев. Систематические ошибки фиксированы для данной машинки (у каждой позиции, рычага или значения — своя), случайные меняются от удара к удару. Модель (\code{tools/analysis.py}) разыгрывает 400 машинок по 400 пар «буква + знак» и складывает ошибки.

\begin{center}\small
\begin{tabularx}{\linewidth}{@{}lXlrr@{}}
\toprule
Исполнение & Звено & Род & $\sigma$, мм & вклад в 95\,\% \\
\midrule
\tabinput{gen/tab_assume.tex}
\bottomrule
\end{tabularx}
\end{center}

\noindent{\small Значения $\sigma$ — допущения, а не замеры, кроме люфта рычага: зазор в прорези сегмента около 0,05~мм на сторону держит рычаг ровно, а больше 0,15~мм уже приводит к столкновениям соседних рычагов при быстром наборе. «Вклад» — насколько уменьшится 95-й процентиль, если звено сделать идеальным.}

\begin{center}\small
\begin{tabular}{@{}lrrrr@{}}
\toprule
Исполнение & медиана, мм & 95\,\%, мм & 99,7\,\%, мм & доля знаков $> 0{,}1$ мм \\
\midrule
\tabinput{gen/tab_mc.tex}
\bottomrule
\end{tabular}
\end{center}

\paragraph{Вывод.} При одинаковых допусках на изготовление литер и подачу челнок даёт посадку знаков вдвое точнее рычажного исполнения: смещение больше 0,1~мм — у полупроцента знаков против одиннадцати. Наибольшие вклады у челнока — штифты позиций и люфт подачи; именно их и стоит ужесточать в первую очередь. У рычажного исполнения доминирует выравнивание рычагов, а его в эксплуатации приходится поддерживать по каждому рычагу отдельно.

%=====================================================================
\section{Регулировка}

Порядок важен: каждая следующая регулировка опирается на предыдущие.

\begin{enumerate}
  \item \textbf{Подача.} Набрать 64 литеры tinco ($w = 10$) подряд; длина должна быть $640\,u$ = 169,3~мм. Отклонение указывает на ошибку передачи от колеса спуска к рейке.
  \item \textbf{Штифты.} Для каждой позиции ударить литеру и проверить её положение относительно соседних на контрольной строке; штифты — заводская регулировка.
  \item \textbf{Нулевой упор ползуна.} Набрать «telco + точка сверху» и «ando + точка сверху» несколько раз; регулировать нулевой упор, пока точка не встанет на ось у обеих букв. Это единственная регулировка, переносящая посадку всех знаков.
  \item \textbf{Фиксаторы секторов.} Контрольный текст (\S\ref{sec:tests}) перебирает все значения $b$ и $c$; при правильной работе каждая точка встаёт на свою ось.
  \item \textbf{Молоток.} Сила пружины — по плотности оттиска самой широкой литеры (ando) и самой мелкой техты (точка); у правильно настроенной машинки обе пропечатываются без провалов.
\end{enumerate}

%=====================================================================
\section{Испытания}
\label{sec:tests}

\subsection{Контрольный текст}

Скрипт подбирает наименьший набор литер, в котором встречаются все значения $w$, $b$ и $c$, и после каждой литеры ставит точку сверху и точку снизу. Получается 13 литер; на оттиске видно, попала ли каждая точка на свою ось.

\begin{figure}[h]
\centering
\setlength\tgem{10mm}
\begin{tikzpicture}[x=\tgu,y=\tgu]
\input{gen/fig_test.tex}
\end{tikzpicture}
\caption{Контрольный текст, рассчитанный по модели ТПМ-2. Красная линия — ось техты, синяя — нижняя ось. На правильной машинке верхняя точка стоит на красной линии, нижняя — на синей (или на красной, если $d = 0$).}
\label{fig:test}
\end{figure}

{\small\setlength\tgem{5mm}
\begin{longtable}{@{}cl rrr@{}}
\toprule
& Литера & $w$ & $b$ & $c$ \\
\midrule\endhead
\tabinput{gen/tab_test.tex}
\bottomrule
\end{longtable}}

\subsection{Программа}

\begin{center}\small
\begin{tabularx}{\linewidth}{@{}>{\raggedright}p{38mm}>{\raggedright}X>{\raggedright\arraybackslash}p{40mm}@{}}
\toprule
Испытание & Как & Критерий (ориентир) \\
\midrule
подача & 64 литеры tinco подряд & длина 169,3~мм $\pm$ 0,3~мм \\
посадка знаков & контрольный текст, лупа с сеткой & смещение точки от оси $\le$ 0,1~мм \\
повторяемость & контрольный текст 20 раз & разброс одной точки $\le$ 0,05~мм \\
коды & контрольный текст: каждая точка на своей оси & ни одного промаха \\
оттиск & ando и точка на одной строке & равная плотность, без провалов \\
блокировки & недожатая клавиша; две клавиши одновременно & машинка не печатает \\
ресурс & 100 страниц контрольного текста & без ухудшения посадки \\
\bottomrule
\end{tabularx}
\end{center}

%=====================================================================
\section{Спецификация узлов}
\label{sec:spec}

\begin{center}\small
\begin{tabularx}{\linewidth}{@{}>{\raggedright}p{50mm}r>{\raggedright\arraybackslash}X@{}}
\toprule
Узел & Кол-во & Замечание \\
\midrule
клавиши с рычагами & 44 & плюс служебные: регистр, два пробела, забой, «назад на литеру», $\pm u$ \\
кодовые пластины & 44 & по два яруса; до 4 выступов на ярус (R, W, B, C или О, N) \\
мерные планки & 6 & R, W, B, C, О, N; рамка подъёма для регистра \\
штифты башенки & 22 & по одному на позицию; тяги от двух клавиш на штифт \\
подъёмы наковальни & 2 & на ряд (передние клавиши) и на два (регистр) \\
челнок & 1 & 4 ряда $\times$ 22, лёгкий металл; сменный \\
наковальня с башенкой & 1 & на U-образной раме \\
молоток с пружиной & 1 & завод от возврата каретки \\
U-образная рама и ползун & 1 & ход 2,4~мм; нулевой упор с регулировкой; пружина возврата \\
сектора памяти & 3 & B, C, W; фиксаторы 0--9 (W: 0--16), защёлки \\
переставной упор спуска & 1 & 4--16 зубьев \\
колесо спуска, шестерня, рейка & 1 & зуб колеса = 1\,u = 0,265~мм хода каретки \\
каретка с пружинным двигателем & 1 & лист по схеме Хаммонда \\
блокировки & 2 & полный ход, очерёдность \\
\bottomrule
\end{tabularx}
\end{center}

%=====================================================================
\section{Этапы макетирования}
\label{sec:stages}

\begin{enumerate}[label=Э\arabic*.]
  \item \textbf{Литеры и кегль.} Моноширинная переделка серийной машинки (или печать литер на бумаге по модели) — проверить, читаются ли тенгвар и техты при em = 1/6''. Решает судьбу кегля до всякой механики.
  \item \textbf{Ползун и сектора.} Стенд на плите: ползун с нулевым упором, три сектора с фиксаторами, мерные планки и одна клавиша с перебираемыми пластинами. Замерить повторяемость упоров, допуск зазоров и силы.
  \item \textbf{Спуск.} Колесо спуска с переставным упором на 4--16 зубьев и каретка: проверить подачу по \S\,«Регулировка», п.~1.
  \item \textbf{Челнок.} Башенка на 22 штифта, челнок на 4 ряда, два подъёма наковальни, молоток: проверить время поворота и оттиск.
  \item \textbf{Сборка.} Всё вместе, блокировки, программа испытаний.
\end{enumerate}

Если удастся найти живой Хаммонд или Varityper, этапы Э3--Э4 сокращаются: башенка, молоток и каретка берутся готовыми, новыми становятся челнок, ползун и кодовая часть.

%=====================================================================
\section{Риски и открытые вопросы}

\begingroup\sloppy
\begin{enumerate}[label=Р\arabic*.]
  \item \textbf{Мелкий кегль.} Точка-техта около 0,5~мм, дуга буквы 1,7~мм. Проверяется на этапе Э1; если не годится — кегль 1/5'' (челнок вдвое тяжелее по инерции, см.~\S\ref{sec:size}).
  \item \textbf{Зазоры кода.} Требование $\sigma \le 0{,}05$~мм — жёсткое для штампованных пластин; возможно, понадобятся регулируемые выступы.
  \item \textbf{Ползун под молотком.} Рама с молотком не должна перекашиваться при отводе и сдвигаться при ударе; нужна проверка жёсткости.
  \item \textbf{Завод молотка.} На строку ударов больше, чем на обычной машинке; завода пружины должно хватать на самую длинную строку.
  \item \textbf{Скорость.} Хаммонд и блокировка очерёдности делают машинку медленной.
  \item \textbf{Заправка листа.} Молоток за листом — заправка неудобнее, чем на рычажной машинке.
  \item \textbf{Память на одну литеру.} Знаки ставятся только к последней букве.
\end{enumerate}
\endgroup

\subsection*{Что нужно для перехода к рабочим чертежам}

\begin{itemize}
  \item решение, строится ли машинка физически, и на какой основе (с нуля или на Хаммонде / Varityper);
  \item замеры живого Хаммонда или Varityper: радиус наковальни, ход клавиши, силы пружин, устройство башенки — они заменят допущения \S\ref{sec:size} и \S\ref{sec:forces};
  \item результат этапа Э1 по кеглю;
  \item рисунок литер для машинного шрифта (Alcarin — книжный шрифт с контрастом; для челнока нужен равномерный штрих), с сохранением метрик.
\end{itemize}

%=====================================================================
\appendix
\section{Коды литер}

$w$ — подача, $a$ — от края клетки до оси техты (для справки), $b$ — отвод для надстрочных знаков, $d$ — смещение нижней оси, $c = b - d$ — отвод для подстрочных. Всё в $u$.

{\small\setlength\tgem{5.5mm}
\begin{longtable}{@{}cllrrrrr@{}}
\toprule
& Литера & U+ & $w$ & $a$ & $b$ & $d$ & $c$ \\
\midrule\endhead
\tabinput{gen/tab_codes.tex}
\bottomrule
\end{longtable}}

\section{Кодовые пластины}
\label{app:plates}

Для каждой клавиши (по позиции QWERTY) и яруса (н — нижний, в — верхний): литера, род, ряд челнока, номер штифта и зазоры выступов в миллиметрах при $T = 14$~мм, $k = 0{,}6$~мм/u. Прочерк — выступа нет.

{\small\setlength\tgem{5mm}
\begin{longtable}{@{}llcl rr rrrr rr@{}}
\toprule
Кл. & Яр. & & Род & Ряд & Штифт & R & W & B & C & О & N \\
\midrule\endhead
\tabinput{gen/tab_plates.tex}
\bottomrule
\end{longtable}}

\section{Отклонённые варианты}

\begin{description}[style=nextline,leftmargin=0pt]
  \item[Рычажная машинка с салазками знаков (ред. 0.1).] Проще по приводу, но выравнивание держится по 44 рычагам и двум салазкам (по Монте-Карло — вдвое хуже посадка знаков), сила удара зависит от пальца, раскладка ограничена.
  \item[Кегль 1/4'' (ред. 0.2).] Челнок в 5 раз тяжелее по инерции, радиус около 81~мм.
  \item[Без нижней оси.] Ошибка подстрочного знака — $d$: 0,26~мм у основных тенгв, 0,8~мм у hyarmen.
  \item[Без «назад на литеру».] Исправление — до 16 ударов забоя.
  \item[Моноширинная переделка.] Простая и точная, но без пропорциональных ширин; используется как этап Э1.
\end{description}

\section*{Источники}
\begin{itemize}\small
  \item G.\,C.~Mares. History of the Typewriter. London, 1909, с.~143 — \url{https://tekniskkulturarv.dk/book/21c52e37-c94c-4fd7-8f2f-df7031f6b69b?page=151}.
  \item The Virtual Typewriter Museum: Hammond~1 (башенка со штифтами) — \url{https://typewritermuseum.org/collection/brands/index.php?machine=hammond1&cat=ks}.
  \item KBD.news: Hammond~2 (челнок 118°, три ряда, подъём наковальни) — \url{https://kbd.news/Hammond-2-typewriter-2140.html}.
  \item Mind Machine: Hammond Typewriter (взвод молотка возвратом каретки; скорость) — \url{https://mindmachine.co.uk/products/04_Manuf_Hammond_01.html}.
  \item Antique Typewriters: Hammond~1 — \url{https://antiquetypewriters.com/typewriter/hammond-1-typewriter}.
  \item Collectors Weekly: Beautiful Machine: 1881 Hammond — \url{https://collectorsweekly.com/articles/?p=16805}.
  \item US~780102 (Хаммонд: молоток, челнок, U-образная рама) — \url{https://patents.google.com/patent/US780102A/en}.
  \item US~2581200 (Varityper) — \url{https://patents.google.com/patent/US2581200}.
  \item US~2904159 (пропорциональный спуск) — \url{https://patents.google.com/patent/US2904159A/en}.
  \item IBM Selectric Composer (единицы 1/72, 1/84, 1/96'') — \url{https://ibmselectric-typewriter.blogspot.com/2011/12/selectric-composer.html}.
  \item Firgelli: Beater Bar (люфт рычага в прорези, отношение клавишного рычага) — \url{https://www.firgelliauto.com/blogs/mechanisms/beater-bar-typewriter}.
  \item HP Museum: Mechanical adders (комптометр) — \url{https://hpmuseum.org/adddir.htm}.
  \item Остальные источники (Alcarin Tengwar, раскладка Смита, IBM Executive) — в документе ТПМ-1.
\end{itemize}

\end{document}
